\chapter{Gastritis}

\section*{Definition}
Gastritis is inflammation of the gastric mucosa, diagnosed histologically rather than endoscopically or symptomatically. It is classified as acute (erosive/hemorrhagic) or chronic (atrophic or non-atrophic), and by etiologic subtype (autoimmune, \textit{H. pylori}-related, chemical/reactive).

\section*{What Happens in Your Body}
Disruption of the gastric mucosal barrier by infectious agents, drugs, or autoimmune attack allows back-diffusion of acid and pepsin, causing inflammation, erosion, and bleeding. Chronic inflammation may progress to glandular atrophy, intestinal metaplasia, and increased gastric cancer risk.

\section*{Causes}
\begin{itemize}
  \item \textbf{Infectious:} \textit{Helicobacter pylori} (most common cause of chronic gastritis worldwide), \textit{H. heilmannii}, CMV (immunocompromised), syphilis, tuberculosis
  \item \textbf{Medications:} NSAIDs (including low-dose aspirin), corticosteroids, bisphosphonates, iron supplements, potassium chloride
  \item \textbf{Autoimmune:} Autoimmune metaplastic atrophic gastritis (AMA) with anti-parietal cell antibodies targeting intrinsic factor and H+/K+-ATPase
  \item \textbf{Chemical:} Bile reflux (post-gastrectomy, incompetent pylorus), excessive alcohol consumption, radiation
  \item \textbf{Other:} Stress gastropathy (critical illness, burns, trauma, major surgery -- Curling's ulcer), ischemic gastritis, portal hypertensive gastropathy, Crohn disease, sarcoidosis
\end{itemize}

\section*{When to See a Doctor}
See your doctor if:
\begin{itemize}
  \item Epigastric pain with hematemesis (coffee-ground emesis) or melena (black, tarry stools)
  \item Severe epigastric pain radiating to the back (suggests penetrating ulcer or pancreatitis)
  \item Iron-deficiency anemia without an alternative source of blood loss
  \item Unexplained weight loss, early satiety, or dysphagia
  \item NSAID use or heavy alcohol consumption with persistent dyspeptic symptoms
\end{itemize}

\section*{Related Symptoms}
\begin{itemize}
  \item Epigastric pain or burning, nausea, vomiting, early satiety, bloating, belching, anorexia, hematemesis, melena
\end{itemize}
\textbf{Important:} Chronic gastritis is often asymptomatic; symptoms correlate poorly with histologic severity, and many people with advanced atrophic gastritis have no pain or discomfort.

\section*{Self-Care / Home Management}
\begin{itemize}
  \item Avoid NSAIDs, aspirin, alcohol, and smoking -- the most common correctable triggers
  \item Eat smaller, more frequent meals rather than large meals that distend the stomach
  \item Avoid spicy, acidic, or fried foods that may exacerbate mucosal irritation
  \item Acetaminophen for pain instead of NSAIDs; antacids (calcium carbonate, magnesium hydroxide) for symptomatic relief
  \item Elevate the head of bed 6--8 inches if nocturnal reflux is also present
\end{itemize}

\section*{How Your Doctor Will Evaluate You}
\begin{itemize}
  \item EGD with gastric mucosal biopsy is the gold standard for diagnosis and classification
  \item Histopathologic features: neutrophil infiltration (acute), mononuclear cells and glandular atrophy (chronic), intestinal metaplasia (pre-neoplastic)
  \item H. pylori testing: urea breath test, stool antigen test, or biopsy-based (rapid urease test, histology, culture)
  \item Serum gastrin and anti-parietal cell/intrinsic factor antibodies for suspected autoimmune gastritis
  \item CBC for anemia, iron studies for iron deficiency (associated with atrophic gastritis)
\end{itemize}

\section*{Treatment Options}
\begin{itemize}
  \item Proton pump inhibitors (omeprazole 20--40 mg daily, pantoprazole 40 mg daily) for 4--8 weeks for acute erosive gastritis
  \item H. pylori eradication: 14-day triple therapy (PPI, amoxicillin, clarithromycin) or quadruple therapy (PPI, bismuth, metronidazole, tetracycline) based on local resistance patterns
  \item Discontinue NSAIDs or switch to COX-2-selective inhibitor if NSAID-induced gastritis
  \item Vitamin B12 supplementation (parenteral) for autoimmune gastritis with pernicious anemia
  \item Surveillance endoscopy every 3 years for extensive atrophic gastritis or intestinal metaplasia
\end{itemize}

\section*{References}
\begin{thebibliography}{99}
\bibitem{gastritis_symptoms:ref1} Sipponen P, Maaroos HI. ``Chronic Gastritis.'' Scand J Gastroenterol. 2015;50(6):657--667.
\bibitem{gastritis_symptoms:ref2} Sugano K, Tack J, Kuipers EJ, et al. ``Kyoto Global Consensus Report on Helicobacter pylori Gastritis.'' Gut. 2015;64(9):1353--1367.
\bibitem{gastritis_symptoms:ref3} Feldman M, Friedman LS, Brandt LJ. ``Sleisenger and Fordtran's Gastrointestinal and Liver Disease.'' 11th ed. Elsevier; 2021.
\bibitem{gastritis_symptoms:ref4} Dixon MF, Genta RM, Yardley JH, et al. ``Classification and Grading of Gastritis: The Updated Sydney System.'' Am J Surg Pathol. 1996;20(10):1161--1181.
\bibitem{gastritis_symptoms:ref5} McColl KEL. ``Helicobacter pylori Infection.'' N Engl J Med. 2010;362(17):1597--1604.
\end{thebibliography}
